% 本文件由 LaTeX 排版 Agent 根据有效 translation.json 自动生成。
% author=John Chrysostom; work_id=1916; title=致奥林比亚书

\chapter{致奥林比亚书}

% structure: p0001 source=h1 target=omitted reason=page-title-already-rendered-by-parent

\EditorialNote{女执事奥林比亚（Olympias）是金口若望（Chrysostom）现存书信中十七封书信的收信人，是他最杰出的女性友人。她出身于异教高官家庭，约生于368年。其父塞琉古（Seleucus）是帝国伯爵，在她幼年时去世，她由舅父普罗科皮乌斯（Procopius）监护，后者是虔诚的基督徒，也是纳西盎的格列高利（Gregory of Nazianzus）之友。格列高利对她极为关心，在书信中称她为\Quote{我的奥林比亚}，并乐于被她称为\Quote{父亲}。她的女教师特奥多西亚（Theodosia），即伊科尼翁的圣安菲洛基乌斯（St. Amphilochius of Iconium）的姐妹，是格列高利劝勉她效法的妇女典范，视为基督徒美德的极致。这位孤女容貌美丽，且继承了大笔财产，故追求者众多。384年，十六岁时她与出身高贵、品行端正的年轻人内布里迪乌斯（Nebridius）结婚。然而婚姻并不幸福，或许正因如此，加上婚后约两年丈夫去世，奥林比亚从中看到天主的暗示，不应再使自己缠累于婚姻生活的世俗忧虑。皇帝特奥多西乌斯（Theodosius）想将她嫁给一位青年西班牙人埃尔皮迪乌斯（Elpidius），他是皇帝的亲属；因她拒绝而恼怒，皇帝下令没收其财产直至她年满三十岁，除非她同意这门亲事。但奥林比亚坚定不移，在一封措辞庄重而讽刺的信中感谢皇帝解除了她的重担：\Quote{除非他下令将她的财富施舍给教会和穷人，否则他不可能赐予她更大的祝福。}特奥多西乌斯意识到其严厉法令无用，或许也后悔不公，遂取消之，让她安享其财产。从此，她的时间和财富都奉献于宗教服务。她照顾病人和穷人的需要，并以慷慨捐赠支持希腊、小亚细亚和叙利亚的教会工作，不仅捐赠金钱，还有土地，以致于那位可称为施舍伟大宣讲者的金口若望也警告她不要不加辨别地慷慨，提醒她说，既然财富是天主托付给她的信托，她应当谨慎管理。这一有益的建议使许多贪财的主教和神职人员对她怀有敌意，这些人曾从她的礼物中获利，或希望从中获利。而奥林比亚则以许多细微的女性关怀回报总主教的灵性关怀，尤其是确保他获得健康的食物，并不要因过于严格的禁欲而过度消耗他孱弱的体质。她自己却践行最严厉的苦修，放弃沐浴的奢侈，只穿旧而粗糙的衣服，并在饮食和睡眠方面严格限制自己。404年，金口若望因敌人阴谋被逐出君士坦丁堡后，奥林比亚因所有追随者所受的迫害而遭受了许多痛苦。她被指控参与了他离开后立即爆发的、烧毁主教座堂和元老院的大火。她在总督面前表现出勇敢的态度，总督试图恐吓她认罪，或诱使她承认阿尔萨西乌斯（Arsacius）——此人通过皇帝的专断权力被强塞进主教席位——但她毫不畏惧，这引起了普遍的钦佩。她坚韧的消息给这位流亡的总主教在肉体痛苦和精神困扰中带来了巨大的安慰。她是否被逐出君士坦丁堡还是自愿退隐，尚不确定；关于她余生的详细情况，我们也没有确切的信息。}\par

% structure: p0003 source=h2 target=section reason=relative-heading-rank; base=h2

\section{致我夫人}

最可敬并蒙天主垂爱的女执事\Person{奥林比亚}，我，主教\Person{若望}，在主内问安。\par

1. 现在，让我来疗愈你沮丧的创伤，驱散那在你周围聚集这忧虑之云的思绪。因为是什么扰乱了你的心？你为何忧伤沮丧？是因为那降临在教会身上的猛烈黑风暴，将万物笼罩在如同无月之夜的黑暗中，且每日加剧，努力制造灾难性的海难，增加世界的毁灭？我与你知道这一切；无人能否认，若你愿意，我将为现在所发生之事勾勒一幅画面，好更清晰地向你呈现这悲剧。我们看到大海从最深处翻腾，有些水手漂浮在波涛上死去，另一些被吞噬，船只的木板碎裂，船帆撕成碎片，桅杆折断，船桨从水手手中脱落，舵手坐在甲板上，双手抱膝而不握舵，以尖锐的呼喊和痛苦的哀叹哭诉他们处境的绝望；天空与大海都不清晰可见，一切是深沉而不可穿透的黑暗，以至无人能看见邻人，巨浪咆哮猛烈，海怪从四面攻击船员。\par

但我要如何继续追求那不可企及的？因为无论我试图用什么形象来描绘我们当前的邪恶，言语都退缩不前，无法成功。然而，即使目睹这些灾难，我也没有放弃对更好事物的希望，因为我考虑到这一切中的舵手是谁——他不是靠技巧战胜风暴，而是用他的杖平息怒涛。但若祂没有在起初就迅速成就此事，这是祂的习惯：祂不在起初就平息这些可怕的灾难，而是当它们加剧，达到极点，大多数人陷入绝望时，祂才以超乎一切期望的方式奇妙地行事，以此显明祂自己的大能，并训练那些承受这些灾难之人的忍耐。因此，不要沮丧。因为只有一件事，\Person{奥林比亚}，是真正可怕的，只有一个真正的考验，那就是罪；而我从未停止反复强调这一主题。至于其他一切，阴谋、敌意、欺诈、诽谤、侮辱、控告、没收、流放、敌人的利剑、深海的危险、全世界的战争，或任何你愿意列举的事，它们都不过是无聊的故事。因为无论这些事物的性质如何，它们都是短暂易逝的，在可朽的身体中运作，而不伤害警醒的灵魂。因此，蒙福的\Person{保禄}，渴望证明与此生相关的喜乐和忧愁的微不足道，他用一句话宣告了全部真理，说：\Quote{看得见的事物是暂时的。}\BibleRef{格后 4:18} 那么，为何你惧怕那些如同河水一样流逝的暂时事物呢？因为现世事物的本质就是如此，无论是令人愉悦还是痛苦。另一位先知将人的一切繁荣不仅比作草，而且比作更脆弱的材料，将其整体描述为\Quote{如草上的花。}因为他没有指出其中任何一部分，如仅财富、仅奢侈、仅权力或仅荣誉；而是将所有在人眼中视为辉煌的事物统称为荣华，然后说：\Quote{凡人的一切荣华，如同草上的花。}\BibleRef{依 40:6}\par

2. 然而，你会说，逆境是可怕的事，难以忍受。但再看看它与另一个图像的比较，然后也学会轻视它。因为敌人施加的辱骂、侮辱、责备、嘲笑以及他们的阴谋，被比作一件破旧的衣服和虫蛀的羊毛，当天主说：\BibleQuote{你们不要怕人的辱骂，也不要因他们的毁谤而惊惶，因为他们必像衣服一样变旧，像虫蛀的羊毛一样被消灭。}\BibleRef{依 50:7} 所以，不要让这些正在发生的事困扰你，而是停止寻求这个或那个人的帮助，不要追逐影子（因为人的联盟就是如此），你要不断呼求你所事奉的耶稣，只需他垂首；转瞬之间，所有这些灾祸都将消散。但如果你已经呼求了他，而它们仍未消散，这就是天主的行事方式（因为我要恢复我之前的论点）：他不一开始就消除灾祸，而是等到它们发展到顶点，等到敌人的恶意几乎每一种形式都已满足时，然后他突然将一切转化为安宁，并引导它们走向意想不到的结局。因为他不仅能将我们所期望和希望的许多事转化为善，而且能转化更多，甚至无限多。因此，保禄也说：\BibleQuote{愿那能远超我们一切所求所想的天主，赐予我们丰盛的恩宠。}\BibleRef{弗 3:20} 例如，他难道不能从一开始就阻止那三个青年陷入试探吗？但他没有这样做，从而给他们带来了极大的痛苦。因此，他容许他们被交到外邦人手中，火炉被烧到无法测量的高度，国王的愤怒比火炉更猛烈，手脚被严严捆住，他们自己被投入火中；然后，当所有看见的人都对他们的得救绝望时，天主的奇迹大能，这位至高工匠，突然超乎一切希望地显现并大放光芒。因为火被束缚了，被捆绑者得到了释放；火炉成为了祈祷的圣殿，一个充满泉源和露水的地方，比王宫更高贵，他们头上的头发甚至胜过了那吞噬一切的火焰，这火焰甚至能战胜铁和石头，并征服每一种物质。在那里，这些圣人们举行了一场庄严的普世赞颂之歌，邀请每一种受造物加入这奇妙的旋律；他们向天主献上感恩的赞美诗，因为他们曾被捆绑，也曾被焚烧，至少就敌人的恶意所能及的程度而言；他们曾是祖国的流放者，失去自由的俘虏，从城市和家园流亡的漂泊者，在陌生和野蛮的土地上寄居；因为这一切都出自一颗感恩的心。当敌人的恶意计谋达到顶点（在他们死后还能做什么呢？），英雄们的劳苦完成，胜利的花环编织完成，赏报准备好，他们的荣誉再无缺欠时，最后他们的灾祸才结束；那点燃火炉并将他们交付那巨大惩罚的人，自己却成为了这些圣英雄的赞颂者，以及天主奇迹的宣告者，并在地四极发出了充满虔诚赞美的书信，记录了所发生的事，成为了施行奇迹的天主所行之事的忠实宣告者。因为他曾是敌人和反对者，所以他写下的东西在敌人看来也是无可怀疑的。\par

3. 你看到属于天主的丰富资源了吗？祂的智慧，祂非凡的能力，祂的慈爱与眷顾？因此不要沮丧或烦恼，而要不断为一切事感谢天主，赞美祂，呼求祂；恳求并祈祷；即使无数骚乱与麻烦临到你，即使暴风雨在你眼前掀起，也不让任何事扰乱你。因为我们的主不会因困难而困惑，即使一切事物都陷入毁灭的极端。因为祂能够扶起跌倒的人，转化误入歧途的人，扶正被诱捕的人，释放被无数罪恶重压的人，并使他们成义，使死人复活，使腐朽之物恢复光彩，使陈旧之物焕然一新。因为若祂能使不存在的事物成为存在，并赋予那无处显明的事物存在，那么祂岂不是更能纠正已经存在的事物？但你会说有许多人丧亡，许多人陷入罗网。这样的事确实常常发生，然而后来他们都得到了适当的纠正，除了一些即使在环境改变后仍处于不可救药状态的人。你为什么因某人被驱逐、某人被取代而烦恼不安？基督被钉十字架，却要求释放强盗巴辣巴，堕落的民众大声呼喊要保存凶手，而非救主和恩主。你想那时有多少人因此跌倒？多少人丧亡？但我必须把论证再往前推。那位被钉十字架者，不是一出生就成了流浪者和逃亡者吗？祂不是从摇篮起就与全家一起迁往异乡，长途跋涉到蛮荒之地吗？这次迁移引发了血流成河、残酷的谋杀和屠杀，所有幼童就像列阵的士兵一样被砍碎，从母亲怀中夺走的婴儿被交付死亡，甚至当乳汁还在他们喉咙里时，刀剑就刺穿了他们的脖颈。还有什么比这悲剧更令人痛苦？这些事是由那企图消灭耶稣的人所做的，然而含忍的天主容忍了这导致如此多流血的悲剧残忍，并忍住没有阻止它，尽管祂有能力，为了某种不可测知的明智目的而显示祂的含忍。当耶稣从异乡回来并长大时，战争在祂周围重新燃起。首先若翰的门徒嫉妒祂，试图毁谤祂，虽然若翰自己对祂恭敬有加，他们说：\BibleQuote{同你在约但河外的那位，你看，祂正施洗，众人都到祂那里去了。}\BibleRef{若 3:26}因为这些话出自那些已经恼怒、被恶意激动、并被此情绪吞噬的人。同样，说这些话的一个门徒曾与某个犹太人争论，并就净化之事挑起争论，比较一种洗礼与另一种洗礼，即若翰的洗礼与基督门徒的洗礼。因为经上记载说：\Quote{若翰的门徒与一个犹太人关于净化发生了一场争论。}当祂开始行奇迹时，有多少毁谤者！有些人称祂为撒玛黎雅人和附魔者，说：\BibleQuote{你是撒玛黎雅人，又附了魔。}\BibleRef{若 8:48}另一些人称祂为\Quote{骗子}，说：\BibleQuote{这人不是从天主来的，而是欺骗民众。}\BibleRef{若 7:12}还有一些人称祂为\Quote{巫士}，说：\BibleQuote{祂是靠着魔王贝耳则步驱魔。}\BibleRef{玛 9:34}他们不断说这些话反对祂，称祂为天主的敌人、贪食贪婪的人、醉汉、罪人与恶人的朋友。因为祂说：\BibleQuote{人子来了，也吃也喝，他们却说：看！一个贪吃好酒的人，一个税吏和罪人的朋友。}\BibleRef{路 7:34}当祂与妓女交谈时，他们称祂为假先知；有人说：\BibleQuote{祂若是先知，必定知道这触摸祂的女人是谁。}\BibleRef{路 7:39}事实上，他们每天磨尖牙齿反对祂。不仅犹太人如此反对祂，连那些被称为祂兄弟的人也不真诚地依附祂，甚至从祂自己的家中也燃起了反对。请看至少他们自己也何等堕落，福音作者附加评论说：\BibleQuote{原来祂的兄弟也不相信祂。}\BibleRef{若 7:5}\par

4. 但既然你想起许多跌倒偏离的人，你以为在祂被钉十字架时，有多少门徒跌倒呢？一个背叛了祂，其余的都逃跑了，一个否认了祂，当所有人都离弃祂时，祂被绑着带走，无人陪伴。\newline{}那么，你以为那些不久前看见祂行奇迹、复活死人、洁净痲疯病人、驱逐魔鬼、增加饼、并行了各种奇事的人，在那个时刻有多少人跌倒呢？当他们看见祂被带走并捆绑，被普通士兵包围，犹太司祭们喧哗吵闹着跟随，祂独自一人被敌人围困，叛徒站在一旁为他的行为得意洋洋？\newline{}你以为当祂被鞭打时，效果如何？很可能有一大群人聚集。因为那是一个显赫的节日，把所有人都聚在一起，这罪行之剧在京城上演，且就在正午。你以为当时在场的人有多少人跌倒，当他们看见祂被捆绑、鞭打、血流如注、在总督审判台前受审，而没有一个门徒站立身边？\newline{}当祂遭受那多种多样接连重复的嘲笑，当他们用荆棘给祂加冕，然后给祂穿上华丽长袍，再把一根芦苇放在祂手里，然后跪下朝拜祂，用尽各种下流和嘲笑时，效果又如何？\newline{}你以为当祂们打祂的巴掌并说：\BibleQuote{对我们说预言吧，基督！打你的是谁？}\BibleRef{玛 26:28}时，有多少人跌倒、多少人困惑、多少人迷茫？当他们把祂带来带去，整天在犹太会众中间嘲笑、辱骂、下流和讥讽？当大司祭的仆人打了祂一巴掌？当士兵分祂的衣服，当祂被带上十字架，背上带着鞭打的伤痕，被钉在木头上时，你以为有多少人跌倒？\newline{}因为即使那时，那些野蛮的野兽也没有软化，反而比以前更狂暴，悲剧更激烈，嘲笑增加了。因为有些人说：\BibleQuote{啊！你这拆毁圣殿、三日内又建造起来的，}\BibleRef{玛 27:40}有些人说：\BibleQuote{祂救了别人，却救不了自己。}\BibleRef{玛 27:42}\par

还有人说：\BibleQuote{如果你是\Person{天主之子}，就从十字架上下来，我们就相信你。}\BibleRef{玛 27:40}\par

再者，当他们用海绵递上苦胆和醋来侮辱祂时，你以为有多少人跌倒？或者当强盗辱骂祂时？或者如我已经说过的，当他们做出那可怕荒谬的断言，认为那个强盗和破门而入者、背负杀人罪的凶手比耶稣更应被释放，并从审判官那里获得许可作选择，反而偏好巴拉巴，不仅渴望钉死基督，还要让祂蒙受丑名？因为他们认为通过这些手段，他们能制造一种信念，认为祂比强盗更坏，是如此的罪大恶极，以至于既不能因怜悯，也不能因节日的特权而救祂。因为他们所做的一切都是为了毁谤祂的名声；这也是为什么他们把两个强盗和祂一起钉死的原因。然而，真理没有被遮掩，反而更清楚地照耀出来。\newline{}他们还指控祂篡夺王权说：\Quote{凡自称为王的，就不是凯撒的朋友。}将这项篡夺的指控加于那没有枕头之地的人身上。\newline{}此外，他们还用亵渎的诽谤指控祂。因为大司祭撕开自己的衣服说：\BibleQuote{祂说了亵渎的话；我们何必还需要见证人呢？}\BibleRef{玛 26:65}\newline{}祂的死不是暴力的吗？不是重罪犯的死吗？不是可憎罪犯的死吗？不是最卑劣的那种吗？不是犯下最恶劣罪行、连在地上吸最后一口都不配的人的死吗？\newline{}至于祂的埋葬方式，难道不是出于恩惠成就的吗？因为有人来乞求祂的遗体。因此，连埋葬祂的人也不是祂的朋友，不是祂曾施恩的人，不是祂的门徒，不是那些曾与祂自由而有益地交往的人，因为所有人都逃跑了，所有人都急忙离开了祂。\newline{}祂的敌人因复活而捏造的卑鄙猜疑，他们说：\BibleQuote{祂的门徒来了，把祂偷走了。}\BibleRef{玛 28:13}你以为有多少人跌倒，有多少人一时为此不安？因为那故事在当时流行，虽然是捏造，且是用钱买来的；但尽管如此，它在一些人中仍然站得住脚，在（坟墓的）封印被打破之后，在真理明显显现之后。因为群众不知道复活的预言（这不足为奇），因为就连祂的门徒也不明白；因为我们读：\BibleQuote{他们不知道祂必须从死者中复活。}\BibleRef{若 20:9}\newline{}因此，你以为在那些日子里有多少人跌倒？然而，长久忍耐的天主耐心地承受，按照祂自己不可测的智慧安排一切。\par

5. 然后，那些日子过后，门徒们继续躲在暗中，秘密地生活，成为充满恐惧和战栗的逃亡者，不断地从一处转移到另一处。甚至当五十天后他们开始出现并施行奇迹时，也并未享有完全的安全；但即使在那之后，仍有无数绊脚石使软弱的弟兄跌倒：当他们被鞭打时，当教会受苦时，当他们自己被驱逐、敌人在许多地方占了上风并引起骚乱时。因为当他们因所行的奇迹而获得很大信心时，\Person{斯德望}的死又引发了严重的迫害，驱散了他们所有人，使教会陷入混乱；门徒们再次惊慌、逃亡、受苦。然而教会不断成长，因所行的征兆而繁荣，并因其创始方式而变得显赫。例如，一个门徒被从窗户缒下，从而逃脱了统治者的手；另一些被天使带出监狱，从而摆脱了锁链；还有的被平民和工匠之家收留，当他们被掌权者赶出时；他们受到各方面礼貌的对待：卖紫色布的妇人、制帐幕者和住在城郊海边的制皮匠。此外，他们常常不敢出现在城镇中心；即使他们冒险去了，他们的接待者也不敢。就这样，在交替的考验和考验的间歇中，教会的大厦被建造起来，曾经跌倒的人后来被扶起，迷失的人被带回，毁坏的地方比先前更坚固地建立起来。为此，当\Person{保禄}祈求圣言的宣讲能平稳进行时，富有智慧和谋略的天主并没有顺从祂的门徒；即使多次恳求，祂也不应允，却说：\BibleQuote{我的恩宠为你足够了，因为我的能力在软弱中才全显出来。}\BibleRef{格后 12:9} 如果现在你甚至将好事与痛苦一起计算，你会看到许多事件发生，即使不是正面的征兆和奇迹，也确实类似征兆，并且是天主伟大眷顾和援助的不可言喻的证据。但为了不让你毫不费力地听我说一切，我将此作为你的任务留下：让你精确地计算一切，并将它们与不幸相比较，并通过从事这项好的工作来转移你的心思，远离沮丧；因为你会从这项工作中得到许多安慰。\par

请代我向您所有蒙福的家人说许多亲切的话。愿最可敬和蒙天主眷顾的女士，您继续享有健康和愉快的心情。\par

如果您希望我写长信，请告知我这一点，并请不要欺骗我说您已摆脱一切沮丧并享受休息的时光。因为信件是一种适当的补救，能在您心中产生极大的喜乐，您将不断收到我的信。下次您写信给我时，不要说：\Quote{我从您的信中得到了许多安慰，}因为我自己知道这一点，但要告诉我，您已得到了如我所愿的那样多的安慰，您没有被忧伤困惑，您不是在哭泣中度过时光，而是在平静和喜乐中。\par

% structure: p0015 source=h2 target=omitted reason=duplicate-page-title

不要为我焦虑，也不要因严冬的酷寒、我消化系统的虚弱以及依扫利亚人的侵扰而过度忧虑。因为这里的冬天无非是亚美尼亚惯常的冬季，无需多言；它并未对我造成严重伤害。因为事先我已想出了许多对策来避免由此可能产生的损害：保持炉火不熄，在居住的房间里设置屏风，使用大量毛毯，并始终待在室内。这确实令我烦恼，若非出于益处——只要我留在室内，寒冷便不会过于困扰我；但若我必须稍作外出并接触外界空气，便会遭受不小的损害。因此，我恳求您，亲爱的夫人，并视此为极大的恩惠，请您十分注意恢复身体的健康。因为沮丧导致疾病；当身体疲惫虚弱、长期处于被忽视的状态，缺乏医生的帮助、有益健康的气候以及充足的生活必需品时，请想想由此会带来多大的痛苦加重。为此，我恳求您，亲爱的夫人，要延请多种医术高明的医生，并服用能够改善这些状况的药物。因为几天前，由于大气状况，我出现呕吐倾向，在采用其他疗法之余，我服用了最审慎的女主人\Person{辛克勒提雍}派人送给我的药物，我发现仅用了三天便治愈了我的不适。所以我恳求您自己也使用这种药物，并安排再送一些给我。因为我再次感到不适时，又使用了它，完全治愈了我的病症——它能平息内部深处的炎症，使水分从皮肤排出，产生适度的温暖，注入不小的精力，并能激发食欲；所有这些效果我都在几天内经历了。愿我最尊敬的主人提奥斐洛伯爵被劝勉设法再送一些这种药物给我。不要因我在此过冬而担忧，因为我现在的健康状况比去年舒适和稳定得多；因此，如果您也注意自我照顾，您的情况会更令人满意。现在，如果您说您的疾病是由沮丧引起的，那么您为何再次向我索要书信呢？您并未从中获得任何欢愉的益处，反而深陷沮丧的专制之下，甚至渴望离开这个世界。难道您不知道，怀有感恩之心的人，即便在病痛中也能获得多大的赏报吗？难道我不是常常亲自并通过书信与您谈论这个主题吗？但由于或许事务的压力，或您疾病特有的性质以及您状况的迅速变化，使得您无法持续清晰地将我的话记在心中，那么请再听一次，我试图通过重复同样的咒语来治愈您沮丧的创伤：\Quote{写同样的事，}经上说，\Quote{在我并不麻烦，对你们却是稳妥的。}\BibleRef{斐 3:1}\par

2. 那么我现在要说和写的究竟是什么？\Person{奥林比亚}，没有什么比在苦难中耐心忍受更能为一个人增添光彩的了。这确实是诸德之后，荣冠的圆满；它超越所有其他形式的义德，而这一种忍耐尤其比其余的更光荣。也许我所说的似乎隐晦，让我试着更清楚地解释。那么我所肯定的到底是什么？不是财产的掠夺，即使一个人被剥夺得一无所有；不是失去荣誉，不是被逐出祖国、流放到远方；不是劳苦和艰辛的压力，不是监禁和束缚，不是侮辱、谩骂和嘲笑（尽管你不可认为勇敢地忍受这类事情是一种微不足道的刚毅，正如\Person{耶肋米亚}那位伟大而杰出的先知所证明的，他因这种考验而备受煎熬）；\EditorialNote{耶肋米亚第十五章}但即使是这些，也不是失去儿女，即使他们在一瞬间被从我们身边夺走；不是敌人不断的攻击，也不是任何其他类似的事——不，甚至那些被认为痛苦之首的死亡，虽然可怕可憎，都不如身体的疾病那样令人难以承受。这一点由那位最伟大的忍耐英雄所证明，他在被身体疾病包围时，认为死亡是从折磨他的灾祸中解脱；而当他承受所有其他苦难时，却几乎感觉不到它们，尽管他一次又一次地受到打击，最后是致命的一击。因为那并不是小事，而是他的敌人最恶毒的残忍的证据，对付一个并非初次经历苦难、也不是初次上阵，而是已被频繁的攻击折磨得筋疲力尽的人，给予那致命的一击——他的孩子们的毁灭，而且如此残忍地加诸于他，以至于所有男女老少在同一时刻，在青春年华，以暴力的方式被消灭，他们的死亡如此突然，以至于也成了他们的埋葬。因为他们的父亲既没有看到他们被放在床上，也没有亲吻他们的手，没有听到他们临终的话语，没有触摸他们的手或膝，没有在他们即将死去时合上他们的嘴或眼睛——这些行为对安慰与子女分离的父母来说，并非无足轻重；他也没有送别其中一些人的葬礼，并在回家后找到另一些人以安慰他对逝者的思念。他听说，当他们正靠在宴席的软椅上时——那是一个充满爱意、而非放纵的宴席，一张兄弟情谊的餐桌——他们全都被压垮了；血和酒、杯子和天花板、餐桌、尘土和他孩子们的肢体，全都混在一起。然而，当他听到这些事，以及在此之前其他令人痛苦的事——那些事也同样令人痛苦：羊群和整群的牲畜被毁灭，后者被从天上降下的火烧尽（那场悲剧的恶信使如此说），前者被各种敌人一起夺走，连同牧羊人一起被砍碎——尽管如此，我说，当他在极短的时间内看到这巨大的风暴席卷了他的土地、房屋、牲畜和儿女，看到一波接一波的巨浪，长长的岩石，深深的黑暗，无法忍受的滔天巨浪时，他并没有被绝望所折磨，几乎对那些发生的事情无动于衷，除非他是人、是父亲。但当他被交付给疾病和疮痍时，他就渴望死亡，他就哀叹自己、悲恸不已，这样你就能明白，这种苦难比所有其他苦难都更严重，而这种忍耐的形式是最高超的。魔鬼自己也深知这一点；因为当它在所有这些试探之后，看到那位英雄依然平静、毫不惊慌时，它就冲向了这作为最大较量的一个，说所有的灾祸都是可忍受的，如失去孩子、财产或任何其他东西（因为这就是\BibleQuote{皮换皮}这话的意思\BibleRef{约 2:4}），但致命的一击是当痛苦加诸于人的身体时。因此，当它在这次较量中失败后，就再也没话可说，尽管在之前的场合它进行了最顽强、最无耻的反抗。但在这件事上，它发现自己无法再发明任何无耻的诡计，而是掩面退却了。\par

3. 然而，不要认为约伯渴望死亡，因为他无法忍受自己的痛苦，这就成为您渴望死亡的理由。请考虑他渴望死亡的时刻，以及他处境的景况——当时法律尚未赐下，先知尚未出现，恩宠尚未像后来那样倾注，他也未从任何其他哲学中获益。因为作为对我们比当时的人要求更多、分配给我们更艰难任务的明证，请听基督所说的话：\BibleQuote{除非你们的义德超过经师和法利塞人的义德，你们决不能进入天国。}\BibleRef{玛 5:20} 因此，不要以为现在祈求死亡是无可指责的，而要听从圣保禄的声音，他说：\BibleQuote{离去与基督同在一起，这要好得多；但为你们的缘故，留在肉身内却更为必要。}\BibleRef{斐 1:23} 因为越是增加患难的压力，胜利的冠冕就越多；金子越被加热，就越被炼净；商人在海上航行时间越长，所装载的货物就越多。所以，不要以为如今分配给您的劳苦是轻微的，相反，它比您所承受的一切都要崇高，我指的是身体软弱所带来的劳苦。以拉匝禄为例\EditorialNote{路加福音第十六章}（虽然我可能已常对您提及此事，但这并不妨碍我现在再说），身体的软弱有助于他的救恩；他死后被送到了那个人的怀里，那人拥有一处与所有过路人共享的住所，并因天主的命令而不断迁移家园，还牺牲了自己的儿子，他的独生子，是在极其年老时赐予他的；尽管拉匝禄未做任何这些事，却因欣然忍受了贫穷、软弱和孤独而获得了这福分。因为对于勇敢承受一切的人而言，这善行如此之大，以至于能够释放任何犯下最严重罪恶的人，除去他们最沉重的担子；或者，如果有人是正直公义的人，这便成为最大的信赖的额外基础。因为对于义人来说，这是明亮的胜利花冠，光辉远超过太阳的光辉；对于犯罪的人，这是最大的净化方式。为此，保禄将那位犯乱伦婚姻的人交给\BibleQuote{肉身的毁灭}\BibleRef{格前 5:5}，以此净化他。因为证明所行的确实净除了如此大的污点，请听他的话：\BibleQuote{使他灵魂在主的日子里得以得救。}\BibleRef{格前 5:5} 当他指责别人犯下另一项非常可怕的罪，即不相称地分享圣桌及其隐秘奥迹，并说这样的人将成为\BibleQuote{对主的身体和血负有罪责}\BibleRef{格前 11:27}——然后观察他如何说他们也从那个严重的污点中得到净化——\BibleQuote{因此你们中间有许多软弱和患病的。}\BibleRef{格前 11:30} 然后又为了证明他们不会局限于这种惩罚状态，但会从中获得一些益处，即摆脱该罪应受的惩罚，他补充说：\BibleQuote{因为我们若审断自己，就不会受审断。但我们受审断时，是被主惩戒，免得我们和世界一同被定罪。}\BibleRef{格前 11:31} 此外，从约伯的例证可以清楚看出，那些生活极其义德的人从这样的惩戒中获益良多：约伯在受苦之后比之前更显赫；还有弟茂德的例子：他虽然是这样一位好人，受托如此重要的职务，并与保禄一起周游世界，不是两三天、十天、二十天或一百天，而是连续多日患病，身体极为衰弱。保禄在说下面的话时表明这一点：\BibleQuote{为了你的胃和你的屡次虚弱，略微用点酒吧。}\BibleRef{弟前 5:23} 那位复活死者的人并没有治愈这个人的虚弱，而是把他留在他疾病的熔炉中，好让他从中获得极其丰盛的信赖。因为保禄本人从他导师那里获得的教训，以及从祂所受的训练，他都传授给了他的门徒。因为虽然他并未遭受身体的软弱，但却受到同样严厉的考验的打击，这些考验带来了许多身体的痛苦。他说：\BibleQuote{赐给我一根刺在肉身中，一个撒殚的使者来击打我}——指打击、绳索、锁链、监禁、拖曳、虐待以及公刑吏鞭笞的折磨。因此，由于无法忍受这些事情给身体带来的痛苦，\BibleQuote{为此我三次祈求主，使我脱离这根刺。}（这里三次指许多次）然后，当他没有得到所求时，因已学到考验的益处，他保持沉默，并因所遭遇的事而喜乐。\par

所以，即使您留在家里，卧床不起，也不要认为您的生命是闲散无用的；因为您所受的痛苦比那些被拖拽、虐待和受行刑者折磨的人更为严重，因为您这极度的软弱中有一位常驻的行刑者与您同在。\par

4. 那么，现在不要渴望死亡，也不要忽视疗愈的方法；因为这确实不安全。为此，\Person{保禄}也劝勉\Person{弟茂德}要最大程度地照顾自己。关于疾病，现在说得够多了。但如果与我的分离引起你的沮丧，期望从中得释放。我现在说这些不只是为了鼓励你，而是我确信事情确实如此。因为若非注定要发生，我想我早就离世了，考虑到加给我的试炼。跳过在\Place{君士坦丁堡}发生的一切，我离开后，你可以想象我在那漫长而残酷的旅途中所受的苦，其中大部分足以致死；我抵达此处后，从\Place{库库苏斯}迁移后，以及在\Place{阿拉比苏斯}逗留后所受的苦。然而我度过了这一切，现在身体康健，十分安全，以致所有亚美尼亚人都惊讶，以我如此虚弱单薄的身躯，竟能忍受如此难耐的严寒，甚至还能呼吸，而那些习惯了冬天的人也都深受其苦。但我至今毫发无伤，逃脱了多次袭击我们的强盗之手，尽管日常缺乏生活必需品，又无法沐浴；虽然从我在此逗留以来一直缺乏这种享受，我现在已如此习惯，甚至不再渴望从它得到安慰，反而比以前更健康。无论是气候的严酷、地区的荒凉、物资的匮乏、缺少随从、医生医术不精、不能沐浴、像囚犯一样永被关在一室、无法活动（而我总需要不断活动）、老与烟火接触、怕强盗、常被围攻，还是其他任何这类事，都没有胜过我；相反，我的健康状况比在别处更好，尽管那时我得到极大的照顾和关注。把这一切都考虑进去，请振作起来，摆脱现在压制你的沮丧，不要对自己施加过度和残酷的补赎。我寄给你我最近写的论文，即\Quote{无人能伤害那不自害的人}，以及我现在寄给你的信，也是为同一论点辩护。所以我求你常读它，如果你的健康允许，大声诵读。因为如果你愿意，它会成为你有效的良药。但如果你与我争辩，不努力治愈自己，不从沮丧的阴郁沼泽中振作，尽管你拥有无尽的劝告和劝勉，我就不会轻易同意经常给你写长信，如果你不能从中得到任何喜乐的好处。我如何知道呢？不是单凭你这样说，而是通过实际的证明，因为你最近断言，正是沮丧导致了你的这场病。既然你亲自作了这宣认，除非你摆脱了身体的疾病，否则我不会相信你摆脱了沮丧。因为如果是前者引起你的失调，如你在信中所说，那么显然当那消散时，另一个也将同时被除去，当根被拔起时，枝子也随之枯萎——如果枝子继续开花繁茂，结出异常的果实，我无法相信你已从痛苦的根源中解脱。因此，不要向我展示言语，而是事实；如果你康复了，你将看到再次寄给你的信超过过去的通信篇幅。那么，认为这不是小小的安慰：我活着，身体健康，在这样的环境中我摆脱了疾病和软弱，我知道这对我的敌人是极大的烦恼和困扰。因此，你应当认为这是最大的鼓励，是你安慰的冠冕。不要称你的家庭为荒凉，它因所受的苦难，如今在天上被指定了一个更高的位置。我为隐修士\Person{柏拉纠}感到极度痛苦。因此，想想那些勇敢坚守立场的人该得到多么大的赏报，当那些在如此纪律和忍耐的惯常中度过时日的人被发现容易堕落时。\par

% structure: p0021 source=h2 target=omitted reason=duplicate-page-title

我从死亡的门口起身，给这位明智的夫人写这封信；我非常高兴，你的仆人们在我终于停泊在港口时遇到了我。因为如果他们在我仍在大海上颠簸、经历肉体疾病的残酷波涛时遇到我，我就很难向你的谨慎精神报喜而不报忧，从而欺骗你。因为这个冬天异常严酷，引发了更加痛苦的内脏风暴，在过去的两个月里，我无异于一个死人，甚至更糟。因为我的生命仅足以感知环绕我的恐怖，白昼、黎明和正午对我而言都如同黑夜，我日夜紧卧床榻，尽管想尽办法，仍无法摆脱寒冷的恶果；但我虽然生着火，忍受着极不愉快的烟雾，困在一个房间里，裹着层层被褥，不敢迈出门槛一步，我遭受了极度的痛苦，头痛之后不断呕吐，食欲不振，持续失眠。就这样，我焦躁不安地度过了漫长黑暗的苦难之海。但为了不让你因我细说痛苦而忧伤，我现在已从所有痛苦中解脱出来。因为春天一近，气温稍有变化，我的一切痛苦便自然消失了。然而，我在饮食方面仍需非常小心；因此我只给胃装清淡的食物，以便它能轻易消化。但得知我明智的女主人濒临死亡，这让我颇为忧虑。然而，出于我对你福祉的深情、焦虑和牵挂，甚至在你的信件到达之前，已有许多人从那里来，带来你康复的消息，这让我从忧虑中解脱出来。\par

如今，我极其欣喜且高兴，不仅得知你已摆脱病患，更因你如此勇敢地承受临到你的一切，称它们为\Quote{虚幻的故事}；而更为重要的是，你竟将这名称也用于你身体上的病患，这是坚强精神的明证，充满勇气的果实。因为不仅勇敢地承受不幸，而且实际上对它们无动于衷，忽视它们，并以如此微小的努力用\OriginalTerm{忍耐}{patience}的桂冠编织你的额冠，既不劳苦，也不辛劳，既不感到痛苦，也不使他人痛苦，反而始终欢喜跳跃、手舞足蹈，这确实是最圆满的\OriginalTerm{哲学}{philosophy}的证明。因此，我欢欣踊跃，跳跃欢喜；我沉浸在喜悦的激动中，对当前的孤独及围绕我的其他烦恼无动于衷，因你灵魂的伟大和你所取得的屡次胜利而感到振奋、明亮，且颇为自豪，这不仅为了你自身，也为了那座庞大而人口众多的城市，在那里你如同一座塔、一个港湾、一堵防御之墙，以榜样雄辩的声音说话，并通过你的痛苦教导男女两性随时为这些竞赛脱去(世俗牵挂)，以全部勇气进入竞技场，并欣然忍受这些竞赛所涉及的辛劳。令人惊奇的是，你无需闯入论坛，也无需占据城市的公共中心，而是一直坐在一间小屋子、一间狭窄的房间内，为竞技的战士服务并膏抹他们；同时，当大海如此在你周围翻腾，波涛汹涌，峭壁、礁石、岩石突出之处和凶猛的怪兽出现在四面八方，一切都笼罩在最深的黑暗中时，你扬起忍耐的风帆，极其平静地漂浮，有如正午时分、晴朗的天气、顺风拂送你；不仅如此，你不仅没有被这痛苦的暴风雨吞没，甚至连水花都没有溅到身上；这非常自然；这就是德性作为舵柄的力量。然而，商人和舵手、水手和航海者，当他们看见乌云聚集，或狂风向他们猛扑，或涌浪泡沫翻腾时，便将船只停泊在港内；如果他们偶然在公海上遭遇风暴，便尽力而为，想方设法将船驶向某个锚地、岛屿或海岸。但你，尽管如此无数的风浪同时向你袭来，大海因风暴的猛烈而从深处涌起，有些人被淹没，有些人漂浮在水面上死去，有些人赤裸地漂在木板上，你却投入灾难的汪洋大海中，称这一切为\Quote{虚幻的故事}，在暴风雨中顺风航行；这非常自然；因为舵手即使在该学科中无限明智，仍没有足够的技术来抵御每一种风暴；因此他们常常畏缩，不敢与波浪搏斗。但你拥有的学问超越每一种风暴——那就是哲学灵魂的力量——它比千军万马更强，比武器更有力，比塔楼和壁垒更坚固。因为士兵所拥有的武器、壁垒和塔楼，只对身体的保护有用，而且并非总是如此，也不是在各方面；但有时所有这些资源都挫败了，使逃向它们寻求庇护的人失去保护。但你的能力并不驱逐野蛮人的武器，也不对抗敌人的计谋，或任何那类的攻击和策略，而是践踏了自然的强制力量，平定了它们的暴政，夷平了它们的堡垒。并且，在不断地与\OriginalTerm{魔鬼}{demons}搏斗时，你赢得了无数胜利，却没有受到一次打击，而是屹立在箭雨之中毫发无伤，并将投向你矛返回给投掷者。这就是你技艺的智慧；通过你所受的痛苦，你报复了那些施加痛苦的人；通过你所遭遇的阴谋，你使你的敌人痛苦，拥有他们的恶意作为名声材料的最佳基础。而你，自己深知这些事，并通过经验获得了洞察力，自然称它们为\Quote{虚幻的故事}。因为，请问，你怎能不这样称呼它们呢？你拥有必死的身体，却藐视死亡，仿佛你急于离开异国，返回自己的家乡；长期遭受最严重的病患，却比健康强壮的人更加愉快；不因侮辱而沮丧，也不因荣誉和荣耀而得意——后者对许多人造成了无限的伤害，他们在司铎职务中享有辉煌的生涯，达到极老的年岁和最可敬的白发后，却因此陷入耻辱，成为那些寻欢作乐者共同的取笑对象。但你相反，身为女子，披着脆弱的身体，易受这些严重攻击，不仅自己避免落入这样的境地，还阻止了许多人这样做。他们确实在竞赛中尚未走远，甚至在起跑点和出发点就被推倒了；而你，在围绕远处的转弯柱无数圈之后，在每一条赛道上都赢得了奖品，在多种角力和搏击中扮演了你的角色。这非常自然；因为德性的较量不依赖于年龄或体力，而只依赖于精神和心志。因此，女人被加冕为胜利者，而男人却被推翻；同样，男孩被宣布为征服者，而年长者却蒙羞。确实，总是应当钦佩那些追求德性的人，尤其当有些人在许多人弃之而去时仍坚守它的时候。因此，我亲爱的女士，你值得最高的钦佩，因为如此众多似乎享有最大声誉的男女老幼都已逃遁，全部在世界眼前伏卧，而且这不是在经过猛烈攻击或敌人惊人的集结之后，而是在交战之前就被推翻，在斗争之前就败北；相反，你经历了如此多的战斗和如此庞大的敌人集结，却远未被削弱或被你的逆境数目吓倒，反而更加充满活力，竞赛的增加给你增添了力量。因为对已成就之事的回忆成为欢喜、喜乐和更大热忱的基础。因此，我欢欣踊跃，跳跃欢喜；因为我不停止重复这一点，并随身带着我喜乐的材料；因此，虽然与你的分离使你痛苦，但你从你的成功业绩中拥有这巨大的安慰；因为我也被放逐到如此遥远的地方，从这个原因中获得了不小的喜乐——我指的是你的勇气。\par

% structure: p0024 source=h2 target=omitted reason=duplicate-page-title

你为何哀叹？你为何折磨自己，要求自己承受你的仇敌都无法向你要求的惩罚，就这样把你的灵魂抛给了沮丧的暴政？因为你经由\Person{帕特里丘斯}之手寄给我的信函，向我揭示了你心灵所受的创伤。因此我也非常忧伤和痛苦，因为当你本应力尽一切努力，以驱除沮丧离开你的灵魂为己任时，你却到处收集令人痛苦的想法，甚至（如你所说）无中生有，毫无理由地将自己撕碎，给你自己造成极大的伤害。因为你未能将我从\Place{库库苏斯}迁走就为此忧伤，这又是为何？然而事实上，就你而言，你确实已经迁走了我，因为你为此尽了全力，付出了努力。即便这事并未实际实现，你也不应为此烦恼。因为也许在天主眼中，我被安排跑更长的双倍赛程，好使胜利的花环更加荣耀，这似乎正是美意。那么，你为何因这些事而烦恼？这些事使我的名声远播，而被认为堪当如此远超过我功绩的伟大荣誉，你也应当欢喜跳跃，载歌载舞，在额上戴上花冠才是。是这地方的荒凉使你忧伤吗？然而，还有什么比我在此地居留更愉快呢？我享有宁静、平安、大量的闲暇和良好的身体健康。因为尽管这城镇既无市场也无集市，但这对我来说不算什么。因为一切事物都像从涌流的泉源中一样，丰沛地倾倒在我身上。我发现主教大人和\Person{狄奥斯库洛}大人，持续忙于为我提供休养。而善良的\Person{帕特里丘斯}也会告诉你，就我在此地的居留来说，我过得愉快而欢喜，被关怀包围。但如果你为发生在\Place{凯撒里亚}的事件哀叹，那么这里你的行为又是有失你身份的。因为在那里，胜利的明亮花环也已为我编织，因为所有人都在宣扬并传播对我的赞美，对我遭受的虐待以及随后的驱逐表示惊叹和诧异。然而同时，不要让任何人知道这些事，尽管它们已成为许多闲谈的话题。因为\Person{佩阿尼乌斯}大人向我透露，\Person{法雷特里乌斯}本人的长老们已经抵达现场，他们声称与我共融，与我的反对者没有交流、往来或合作。因此，为避免使他们不安，不要让任何人知道这些事。因为所遭遇的事确实非常痛苦：即使我没有遭受其他任何苦难，在那里发生的事件也足以为我赚取无数的赏报：我所遇到的危险是如此极端。现在我恳求你保守这些秘密，因此我将向你简要述说，不是为了令你忧伤，而是为了让你喜乐。因为我的收益的素材就在于此，我的财富就在于此，我摆脱罪过的方法就在于此——我的行程不断被这类考验包围，而且是从那些完全意想不到的人加诸于我。因为当我快要进入\Place{卡帕多西亚}地区时，刚从那个几乎以死亡威胁我的\Place{加拉提亚}人手中逃脱，路上有许多人遇见我，说：\Quote{法雷特里乌斯大人正在等候你，四处走动，生怕错过与你相见的喜悦，竭尽全力要见你、拥抱你，并向你展示各种亲切的关怀；为此他发动了男女隐修院。}当我听到这些事时，我未曾预料其中任何一件会真的发生，但在自己的心中却形成了完全相反的印象：但我对此未向那些带给我消息的任何人透露只字。\par

2. 当我于某天深夜抵达\Place{凯撒勒雅}时，已是精疲力竭，浑身酸痛，正值高热极度发作，虚弱不堪，痛苦万分；我偶然找到一家位于城郊的旅店，并竭尽全力寻找几位医师，以平息这炽热的高烧；因为那时正是我的间日热最严重的时候。此外，还有旅途的疲惫、劳苦、紧张、完全没有随从、补给困难、缺少医生、辛劳、炎热与失眠的消耗；因此当我进城时，几乎像个死人。那时，全体神职人员、百姓、隐修士、修女、医师都来看望我，我得到了极大的关怀，每个人都给我各种照顾和帮助。然而即便如此，因高热引起的嗜睡使我极度痛苦。终于，疾病逐渐消退减轻。但是\Person{法雷特里乌斯}却始终没有露面；他却等着我离开，我不知道他究竟有何意图。那时，我见我的病情稍有好转，便开始计划我的行程，以便抵达\Place{库库苏斯}，在路上的灾难之后稍事休息。当我正处在这种境况时，突然传来消息：无数\Person{依扫里亚人}正在蹂躏凯撒勒雅地区，烧毁了一个大村庄，并且极为凶暴。护民官听到这个消息后，带着他所有的士兵出发了。因为他们害怕敌人也会袭击城市，所有人都惊恐万分，极度的恐慌中，国家的土地岌岌可危，甚至连老人都承担起城墙的防卫。当事情处于这种状况时，黎明时分，突然有一群暴徒般的隐修士（因为我必须这样称呼他们，以此表明他们的疯狂）冲进我们所在的房子，威胁要放火烧毁它，并要对我们施以最严重的暴力，除非我们离开。无论是依扫里亚人的威胁，还是我自己深受其苦的疾病，或其他任何事情，都不能使他们更理智，他们继续推进，被如此凶猛的愤怒所驱使，甚至连总督的士兵都感到害怕。因为他们不断用拳脚威胁他们，并夸耀说他们曾可耻地殴打过许多总督的士兵。士兵们听了这些话，便到我这里寻求庇护，恳求我，说：\Quote{即使我们落入依扫里亚人手中，也求你救我们脱离这些野兽。}总督听说后，急忙赶到房子，想帮助我。但隐修士们毫不理会他的劝诫，实际上他也无能为力。他察觉到事态的严重，既不敢建议我出去送死，也不敢让我留在家中，因为这些人过于狂暴，于是他派人去找\Person{法雷特里乌斯}，恳求他因我的疾病和迫近的危险宽限几天。但即便如此，也毫无成效；第二天，隐修士们来得比之前更加凶猛，没有一个司铎敢于站在我身边帮助我，反而羞愧得面红耳赤（因为他们说这些事是按\Person{法雷特里乌斯}的指示做的），他们躲藏起来，即使我叫他们也不回应。何必多说呢？尽管如此巨大的恐怖迫在眉睫，死亡几乎必然，热病又压迫着我（因为我尚未从那引起的困扰中得到缓解），我于正午时分扑入轿中，被人抬出去。全体百姓尖声哭嚎，诅咒这事的肇事者，人人哀泣悲叹。但当我出城后，一些神职人员也渐渐出来护送我，边走边哀悼。听到有人说：\Quote{你们要把他带到哪里去送死？}一个我深为亲近的人对我说：\Quote{我恳求你，离开吧；落到依扫里亚人手中，只要你远离我们。因为无论你落在何处，你都会落在安全的地方，只要逃脱我们的手。}善良的\Person{塞琉西亚}，我主上\Person{鲁菲努斯}的慷慨妻子（她是我最关怀备至的朋友），听了并看到这些事，便劝勉并恳求我住在她位于郊区的房子，那里离城约五英里，她派了一些人来护送我，于是我就动身去了那里。\par

3. 但即使在那里，这针对我的阴谋也没有结束。因为\Person{法雷特里乌斯}一得知她所做的事，据她说，便对她发出了许多威胁。但当她将我接进她的郊区别墅时，我对此一无所知；因为她出来迎接我时对我隐瞒了这些事，却向在场的管家透露，并吩咐他为我提供一切可能的休息条件，若有任何隐修士前来攻击，想要侮辱或虐待我，他就应从她其他田庄召集劳工，以此组织力量对抗他们。此外，她恳求我躲进她的房子，那房子有堡垒且坚不可摧，使我得以逃脱主教和隐修士的手。然而，我未能被说服这样做，而是留在别墅里，对之后筹划的计谋一无所知。因为即使那时他们仍不满足于对我停止愤怒，而\Person{法雷特里乌斯}正如那位女士所说，紧紧逼迫她，严厉威胁她，强迫她甚至把我从郊区赶走，以致于在半夜，我一无所知，那位女士无法忍受他的纠缠，在我不知情的情况下宣布野蛮人已到近前，因为她羞于提及自己所受的强迫。于是，在半夜里，司铎\Person{埃维提乌斯}来到我这里，把我从睡梦中唤醒，大声喊道：\Quote{起来，我求您，野蛮人已到我们眼前，他们就在附近。}请想象我听到这话时的处境！然后，当我问他我们必须做什么？我们无法逃进城去，以免遭受比伊苏利亚人将要对我们做的更糟的事，他强迫我出去。那是半夜，一个黑暗、阴沉的夜晚，没有月亮——这一情况加剧了我们的困惑——我们没有同伴，没有助手，因为所有的人都抛弃了我们。然而，在恐惧的压力下，想到即将到来的死亡，我站了起来，尽管在受苦，吩咐人点亮火把。但司铎吩咐熄灭它们，据他说，以免火光吸引野蛮人攻击我们；所以火把被熄灭了。然后，驮着我轿子的骡子跪倒在地，因为道路崎岖、陡峭且多石，我摔了出来，险些丧命；之后我下了轿子，被人拖着步行，被司铎\Person{埃维提乌斯}紧紧抓住（因为他也从他的骡子上下来了），就这样我蹒跚前行，被拉着，或者更确切地说，被手拽着，因为走这样的险路是不可能的，而且在半夜里，在陡峭的山间。想象我的遭遇，被如此灾难包围，又被热症所困，对已制定的计划一无所知，却恐惧野蛮人，因想到会落入他们手中而战栗。难道您不认为，仅这些苦难，即使除此之外没有别的临到我，也足以消除我的许多罪过，并为我提供大量材料以在天主面前获得称赞吗？这一切的原因，至少依我猜测，是因为我一到\Place{凯撒勒雅}，那些在职位上的人、前任代理官、前任总督、前任护民官，甚至全体百姓，每天都来拜访我，对我极为尊重，视我如掌上明珠；我想这些事激怒了\Person{法雷特里乌斯}，那将我逐出\Place{君士坦丁堡}的嫉妒，甚至在这里也不肯放过我。至少这就是我所猜测的，因为我并非肯定声称，只是怀疑事实如此。\par

还有什么事可说关于路上发生的其他事件、恐惧和危险？当我日复一日地回忆它们，并不断铭记于心，我因快乐而兴奋，我欢喜跳跃，如同一个为自己积蓄了巨大宝藏的人；因为这就是我对它们的立场和感受。因此我也恳求阁下为这些事欢喜、快乐、跳跃，并光荣天主，因他以为我配受这样的苦难。我恳求您将这些事保密，不向任何人透露，尽管大部分总督府的士兵自己经历了极度的危险，却能将故事传遍全城。\par

4. 然而，切勿让任何人从您的明智得知此事，反而要制止那些谈论的人。但若您担心我受苦的后果会持续，请确切知道，我已完全摆脱它们，而且我身体状况比在\Place{凯撒勒雅}逗留时更好。您为何害怕寒冷？因为已经为我准备了一所合适的住所，我的主\Person{迪奥斯库鲁斯}安排一切，使我丝毫不觉寒冷。若我从经验的开始推测，现在的气候在我看来具有东方特色，不亚于\Place{安提阿}。如此温暖，如此宜人。但您说\Quote{也许您因我疏忽而生气}，这令我非常伤心，然而我在许多天前就已致函阁下，恳求您不要把我从这地方移走。现在我还有机会考虑，您需要强有力的辩护和许多辛劳才能为这句话做出满意的道歉。但也许您已做了部分道歉，您说\Quote{我通常忙于思考如何增加我的痛苦}。但我反过来认为您说\Quote{我以思考来增加我的悲伤而自豪}是最严重的指控；因为当您应当尽一切可能驱散痛苦时，您却听从魔鬼的意愿，增加您的沮丧和悲伤。难道您不知道沮丧是多么大的邪恶吗？\par

关于\Person{依撒乌黎人}，今后请打消对它们的恐惧：因为他们已返回本国；总督已为此做了所有必要之事；我在此比在\Place{凯撒勒雅}时安全得多。因为今后除了少数几位主教外，我无需惧怕任何人。至于\Person{依撒乌黎人}，不必害怕：他们已经退去，冬季来临时便困守家中，尽管圣神降临节后他们可能出来。你说没有收到我的来信，这是什么意思？我已经寄了三封长信给你：一封由总督府士兵带去，一封由\Person{安多尼}带去，第三封由我的仆人\Person{安纳多略}带去；其中两封是有益的良药，能唤醒任何沮丧或跌倒的人，并引领他进入健康的宁静状态。你收到这些信后，要不断仔细研读，便会领会其力量，体验其治愈之力和益处，并告知我你从中得益匪浅。我还有另一封与这些相似的信已准备好，但此刻不打算寄出，因为我被你所说的话极其恼怒，你说：\Quote{我积累忧思，甚至编造不存在之事}，这是与你不相称的话语，使我羞愧低头。但请阅读我已寄出的那些信，你将不再说这些话，即使你执意沮丧。至少我未曾停止，也不会停止说：罪是唯一真正令人痛苦的事；其他一切不过是尘土与烟雾。坐监与戴锁链有何痛苦？受虐待能带来如此多的收益，为何痛苦？流放或抄没财产为何痛苦？这些只是言辞，没有可怕的现实，是无忧的言辞。你若说死亡，你只是提到自然之债：一件无论如何必须经历的事，即使无人催促；你若说流放，你只是提到更换国度、见识许多城市；你若说抄没财产，你只是提到自由和从忧虑中解脱。\par

5. 不要停止关注\Person{玛路塔}主教，尽你所能，将他从坑中拉上来。因为为了\Place{波斯}的事务我特别需要他。请尽可能从他那里查明他在那里完成了什么，他为何回家，并告诉我你是否已将寄给他的两封书信交给了他。如果他愿意写信给我，我会再给他回信；但如果他不愿意，至少让他告知你的智慧，那里是否发生了更多的事，以及他再去那里是否可能完成什么。为此我很想与他见面。尽管如此，凡你所能的事都要做，并注意完成你这一部分，即使所有人都冲向灭亡。这样你的赏报才会成全。因此务必尽可能与他交朋友。我恳求你切勿忽视我即将说的话，而要留心注意。\Person{塞拉比翁}主教常隐藏在那里的\Person{马尔西}和\Person{哥特}隐修士告诉我，执事\Person{莫杜瓦留斯}带来了消息，说那位杰出的主教\Person{乌尼拉斯}，即我最近祝圣并派往\Place{哥提亚}的，在成就了许多伟大事迹后已安息。那执事带来了\Person{哥特人}的国王的一封信，请求给他们派一位主教。既然我看不到任何其他方法来应对那迫在眉睫的灾难并加以纠正，只有拖延和推迟（因为他们目前无法航行进入\Place{博斯普鲁斯}或那些地区），请安排因冬季而暂时推迟他们：绝对不能忽视此事：因为这事极其重要。因为有两件事如果发生，会特别令我痛苦，愿天主不许：一件是由这些行如此大恶的人任命一位主教，他们无权任命；另一件是不加考虑地任命任何人。你自己知道，他们并不急于选立一个合适的人为主教；如果这事发生，愿上天不许，你知道后果如何。因此要尽力防止这两件事中任何一件发生；但如果\Person{莫杜瓦留斯}能够安静秘密地速来见我，那将大有益处。但如果这不可能，那就做那时可行的事。因为金钱的事例，以及福音中寡妇的实例，也适用于实际事务。因为那个穷寡妇将两个小钱投入银库，超越了所有投入更多者，因为她用尽了自己全部所有。同样，那些尽心尽力致力于眼前工作的人，就他们而言，便完全完成了它，即使毫无结果，他们也会得到成全的赏报。\par

我非常感谢\Person{希拉留斯}主教：因为他写信给我，请求允许返回本国，在那里料理事务，然后再回来。由于他的在场很有帮助（因为他是一位虔诚、坚定不移且热心的人），我鼓励他离开并迅速返回。因此请注意将信迅速安全地送达他手中，不要被搁置一旁：因为他热切恳求我的信，他的在场大有益处。因此务必照料好这些信；如果长老\Person{赫拉狄乌斯}不在场，请确保通过一位有头脑的谨慎之人之手交给我朋友们。\par

% structure: p0033 source=h2 target=omitted reason=duplicate-page-title

您的虔敬所遭遇的，并无任何奇怪或异常之事，而完全是自然而合理的事，即通过不断的考验，您灵魂的筋骨变得更加强韧，您对于奋斗的热忱与精力得以增长，并且您因此获得极大的喜乐。因为患难的本性正是如此——当它抓住一个勇敢而高贵的灵魂时，它往往会产生这样的效果。正如火在炼金时能使金子更见真纯，同样，患难临到纯金般的品格时，也使之更加纯净、更经过考验。为此，\Person{保禄}也说：\Quote{患难生忍耐，忍耐生老练。}\BibleRef{罗 5:3}因此，我也欢喜跳跃，并因思念您的刚毅而从我独处的处境中获得最大的安慰。为此，虽有无数豺狼围困您，众多恶党聚集，我毫不畏惧；我祈祷，愿现有的诱惑被压制，愿新的不再发生，如此便成全了主的训诫，他命我们祈祷，使我们不陷于诱惑；但若这些事情再一次被允许发生，我对您那黄金般的灵魂抱有极大的信心，它因此为自己获取最伟大的财富。因为他们还能用什么方法恐吓您呢？他们胆大妄为，自取灭亡。是夺去财产吗？但我深知，这些被您视为尘土，比污泥更不值钱。或是流放异地，逐出家园？但您懂得如何住在繁华的大城中，如同住在无人之地，将全部时间用于宁静与安息，把世俗的野心践踏在脚下。或是用死亡来威胁？您也早已在想象中不断操练了；倘若他们把您拖去宰杀，他们拖去的不过是一具已经死去的身体。何须多言？没有人能对您做任何这类的事，而您不早已充分使自己经历过了。因为您始终行走在狭窄而艰辛的道路上，您已在这一切事上训练了自己。因此，您在训练过程中操练了这最美好的技艺，如今在竞赛本身中，您更加荣耀地闪耀，不仅丝毫不为所发生的事所扰动，反而兴奋、跳跃、欢舞。您在训练中预尝的竞赛，如今您轻松地承担，尽管是在一具比蛛网更脆弱的女子身体中，却以轻蔑的嘲笑践踏那些对您咬牙切齿的壮汉的狂怒；并且准备好承受比他们为您预备的更坏的事。有福啊，三倍有福，因那将要赢得的胜利冠冕，但更因这竞赛本身。因为这些斗争的本性正是如此：甚至在奖品颁发之前，在争斗之中，它们就已拥有报偿和赏报——就是您如今正享受的喜乐、欢欣、勇气、坚忍、忍耐、那不被俘获和征服、超越一切的能力；那使您对他人的任何恐怖都无动于衷的完美训练，那在患难的巨浪中屹立于磐石之上的力量，以及当四周大海翻腾时乘着顺风平静航行的能力。这些就是患难在今世、在赢得天国之前所赐的赏报。因为我深知，就在此刻，您因喜乐而高昂，不认为自己披戴着身体；若有机会召唤您这样做，您脱去身体会比别人脱去所穿的衣服更为爽快。因此，您要欢喜快乐，为您自己，也为那些死于有福之死的人——不是在床上，不是在家里，而是在监狱、锁链和折磨中；您只应为那些做这些事的人哀哭，为他们悲伤。但既然您也想知道我身体的健康情况，让我告诉您，我已从最近压迫我的虚弱中暂时解脱，现在处于较为舒适的状态；唯一担心的是，冬天再度来临，可能又严重损害我脆弱的消化系统；至于依索利亚人方面，我们现在享有很大的安全。\par

\SourceRecord{Translated by W.R.W. Stephens.；Nicene and Post-Nicene Fathers, First Series；Vol. 9.；Edited by Philip Schaff.；Buffalo, NY: Christian Literature Publishing Co.,；1889.；<http://www.newadvent.org/fathers/1916.htm>.}
